\chapter{视觉模型}
\begin{quote}\small
\textit{本章摘要。}本章围绕图像中的缺陷识别展开：图像分类判断是否存在缺陷，目标检测给出缺陷的位置，图像分割描出缺陷的边界。先以卷积网络和 ResNet 解释如何从像素形成类别判断，再以 ViT 解释图像区域之间的注意力交互。随后以 Faster R-CNN 和 U-Net 分别讲解框级定位与像素级预测，并由它们的局限引出单阶段检测、多尺度特征和实例分割。每个模型都沿输入、计算、输出和训练目标展开，最后回到细小缺陷、标注不足与跨相机验证这些典型问题。
\end{quote}
风险预警告诉维护人员何时需要关注设备，却不能指出缺陷位于哪里。本章转向巡检图像和表面质检：先判断是否存在缺陷，再定位区域，最后估计边界与尺寸。这些任务使用图像级、框级或像素级标签，与未来故障标签不同。拆检后才能获得的照片可以支持离线诊断，却不能直接作为拆检前预警的输入。
图像模型处理的是像素之间的空间关系。同样一组深浅像素，排列成细长的连续区域时可能表示裂纹，散布在表面时却可能只是纹理。模型既要识别这些关系，也要按任务要求保留位置。本章始终使用表面质检作为例子，逐步回答三个问题：图中有没有裂纹，裂纹在哪里，哪些像素属于裂纹。

\section{先确定任务：类别、位置还是边界}
图像分类为整幅图给出类别。例如，标签可以是“正常”“裂纹”“划痕”，模型输出各类得分或概率。只有一个互斥标签的分类任务假定每幅图归入一类。若裂纹与划痕可以同时出现，就需要为各类别分别标注是否存在，并使用多标签输出。分类适合工件筛查，但一个类别概率不能直接提供裂纹坐标。

目标检测为图中每个目标输出类别和矩形框。它适合寻找多个缺陷、统计数量以及指引复检位置。图像分割则为每个像素给出标签，适合描述不规则边界和计算缺陷面积。语义分割只区分像素类别，实例分割还区分属于不同目标的像素。同类的两条裂纹可以共用一个语义标签，却需要不同的实例编号。像素长度和面积要换算为物理尺寸，还需要相机标定或已知尺度，模型输出本身不提供毫米单位。

分类需要图像级标签，检测需要框标注，分割需要像素标注。标签粒度决定了能直接训练什么输出。只有“有无裂纹”的标签时，不能直接计算逐像素分割损失。若任务只要求决定是否复检，可以先做分类。若必须指出复检位置或估计裂纹形状，就要分别建立检测或分割任务。

这些任务通常共享一个提取图像特征的网络主体，称为\textbf{骨干网络}。骨干之后的\textbf{任务头}将特征转成类别、框或掩码。卷积神经网络（Convolutional Neural Network，CNN）与视觉 Transformer（Vision Transformer，ViT）是两种常用的视觉表示结构。下面先通过分类走通从像素到预测的过程，再解释定位任务为何需要改变这个过程。
\begin{figure}[htbp]
\centering
\begin{tikzpicture}[node distance=0.65cm and 0.65cm]
 \node[idi input, minimum width=1.7cm] (image) {工业图像\\$H\times W$};
 \node[idi process, right=of image, minimum width=2.0cm] (backbone) {视觉骨干\\CNN / ViT};
 \node[idi state, right=of backbone, minimum width=2.0cm] (features) {空间特征\\单尺度或多尺度};
 \node[idi output, right=of features, minimum width=1.8cm] (class) {分类\\图像级};
 \node[idi output, below=0.75cm of class, minimum width=1.8cm] (detect) {检测\\框级};
 \node[idi output, above=0.75cm of class, minimum width=1.8cm] (segment) {分割\\像素级};
 \draw[idi arrow] (image)--(backbone); \draw[idi arrow] (backbone)--(features);
 \draw[idi arrow] (features)--(class); \draw[idi arrow] (features)--(detect); \draw[idi arrow] (features)--(segment);
\end{tikzpicture}
\caption{视觉系统先由卷积或 Transformer 骨干提取表示，再按标签粒度输出分类、检测或分割结果。}
\label{fig:vision-output-granularity}
\end{figure}
图\ref{fig:vision-output-granularity}中的箭头表示从图像到预测的计算。骨干可以提供不同层的特征，但原始 ViT 通常输出单一尺度的图像块表示，不能默认所有骨干都自带多尺度输出。分类可以汇总位置，检测与分割则要继续使用空间对应关系。

\section{图像分类：从局部卷积到 ResNet}
\subsection{卷积如何从像素中提取特征}
设彩色图像为$X\in\R^{H\times W\times C_{\rm in}}$，其中$H,W$为像素高度和宽度，$C_{\rm in}=3$时三个输入通道分别表示红、绿、蓝。隐藏层也有通道，但这些通道表示学习到的特征，不再对应颜色。张量可以理解成带多个索引的数表，本章采用高、宽、通道的轴顺序。

若把图像直接展平后送入全连接层，每个位置都可以拥有独立权重。这样做能计算预测，却没有利用“同一种局部形状可能出现在不同位置”这一规律。卷积用一个小窗口读取局部像素，并在整幅图上重复使用同一组权重。不同卷积核可以响应不同的局部变化，训练则决定哪些响应对识别缺陷有用。

\paragraph{一个窗口怎样变成一个输出。}
取单通道窗口$\left[\begin{smallmatrix}1&2\\3&4\end{smallmatrix}\right]$与核$\left[\begin{smallmatrix}1&0\\0&-1\end{smallmatrix}\right]$，偏置为零，则对应输出为$1\times1+2\times0+3\times0-4\times1=-3$。核移到右边后仍使用这四个权重，只更换被读取的像素。若随后使用整流线性单元（Rectified Linear Unit，ReLU），即$\max(0,z)$，这一响应变为零。多个核可学习不同方向的响应。核权重由训练损失调整，示例核仅用于理解计算，不是要求人工为每种裂纹设计模板。

步幅是相邻窗口起点移动的像素数，填充是在图像边缘补值，空洞率是核内相邻采样点的间隔。例如$5\times5\times3$输入经过$3\times3$核、步幅一、无填充、四个输出通道，得到$3\times3\times4$，含$3\times3\times3\times4+4=112$个参数。每个输出通道都汇总三个输入通道，并不是分别对红绿蓝得到三个互不交流的答案。

\paragraph{卷积的索引、参数与等变性。}
设输入$X\in\R^{H\times W\times C_{\rm in}}$，方形核大小$k$，输出通道数$C_{\rm out}$，步幅$s$、填充$p$、空洞率$d_a$。采用从零开始的空间索引，有
\begin{equation}
 Z_{i,j,c}=b_c+\sum_{u,v=0}^{k-1}\sum_{a=1}^{C_{\rm in}}
 K_{u,v,a,c}X_{si-p+d_au,\,sj-p+d_av,\,a},
\end{equation}
越界输入按指定规则填充。这个深度学习中惯称的卷积实际采用互相关索引；若使用数学卷积，只需翻转核。沿高度方向，为使最下端采样不超过填充后的输入，需满足$s i+d_a(k-1)\leq H+2p-1$，因此
\begin{equation}
 H_{\rm out}=\left\lfloor\frac{H+2p-d_a(k-1)-1}{s}\right\rfloor+1,
\end{equation}
宽度同理。参数量为$k^2C_{\rm in}C_{\rm out}+C_{\rm out}$，与图像面积无关；不共享的局部连接则须为每个输出位置另存一组权重。

在无限网格或忽略边界、$s=1$时，令$(T_\delta X)_{i,j}=X_{i-\delta_1,j-\delta_2}$，直接代入求和式得$K*(T_\delta X)=T_\delta(K*X)$。这就是平移等变性，而不是平移后输出保持不变。步幅$s>1$时通常仅保留对$s$整数倍平移的对应关系，边界填充还会引入额外差异。图像分类可以通过汇聚追求不变性，定位任务却应保留位置随输入一起移动的关系。

\subsection{多层卷积如何形成整幅图的判断}
单个小窗口只能看到裂纹的一部分。把卷积与非线性激活连续堆叠后，后层可以组合前层的局部响应，例如把多个方向一致的边缘组合为较长的结构。网络常用步幅大于一的卷积或池化降低空间分辨率，这个过程称为下采样。它减少计算量，但也可能让细裂纹与背景混在一起。

一个深层位置在输入中能够访问的范围称为\textbf{理论感受野}。令$r_l$表示第$l$层的感受野宽度，$j_l$表示该层相邻位置在原图中的间隔。以$r_0=j_0=1$开始，对于核大小$k_l$、步幅$s_l$且不使用空洞采样的层，有
\begin{equation}
 r_l=r_{l-1}+(k_l-1)j_{l-1},\qquad j_l=s_lj_{l-1}.
 \label{eq:vision-receptive-field}
\end{equation}
新增的$k_l-1$个间隔，每个跨越$j_{l-1}$个输入像素。三层$3\times3$卷积的步幅依次为$1,2,1$时，$(r_l,j_l)$依次为$(3,1),(5,2),(9,2)$。第三层每个内点理论上访问九像素宽的区域，相邻输出却只相隔两像素。这个范围不表示每个像素贡献相同，边界处还可能读到填充值。

设最后一层特征为$F\in\R^{H'\times W'\times D}$，其中$D$为通道数。分类可以先对每个通道作全局平均，再进行线性分类：
\begin{equation}
 v_c=\frac{1}{H'W'}\sum_{i=1}^{H'}\sum_{j=1}^{W'}F_{i,j,c},\qquad
 z=Uv+b,\qquad p_k=\frac{e^{z_k}}{\sum_{a=1}^{K}e^{z_a}}.
 \label{eq:vision-classification}
\end{equation}
这里$K$为互斥类别数，$v\in\R^D$为整图表示，$U\in\R^{K\times D}$与$b\in\R^K$为分类头参数，$z_k$是类别$k$的实数得分，$p_k$是模型给出的类别概率。若真实类别为$y$，单幅图的交叉熵损失为$-\log p_y$。反向传播同时调整分类头和卷积核，使特征对区分类别有用。

全局平均将不同位置汇总为一个向量，因此适合回答“有没有裂纹”。它同时舍弃了直接读取位置的接口。检测与分割必须在汇总之前保留空间特征。分类训练得到的特征也未必足以辨认所有细小缺陷，仍需检查图像分辨率和训练标签。

\subsection{ResNet：网络加深以后怎样训练}
增加层数能让模型组合更大范围的局部信息，但普通深层网络可能比浅层网络具有更高的训练误差。若连训练集都拟合得更差，就不能只用过拟合解释。更长的连续变换使优化更困难，深层网络也不容易自行学出“保持已有表示不变”的操作。

残差网络（Residual Network，ResNet）将一个块的输出写为\cite{he2016resnet}
\begin{equation}
 h_{l+1}=h_l+F_l(h_l).
 \label{eq:vision-residual}
\end{equation}
$h_l$是输入特征，$F_l$是由卷积、归一化和激活组成的残差分支。若期望映射为$H(h_l)$，残差分支要学习的就是差值$H(h_l)-h_l$。当已有表示已经合用时，令分支输出接近零即可保留输入。图\ref{fig:vision-residual-block}展示了分支与捷径如何在同一位置相加。

\begin{figure}[htbp]
\centering
\begin{tikzpicture}[node distance=0.8cm]
 \node[idi input] (h) {$h_l$};
 \node[idi process,right=of h,minimum width=3cm] (f) {卷积、归一化、激活\\$F_l(h_l)$};
 \node[idi state,right=of f] (add) {$+$};
 \node[idi output,right=of add] (out) {$h_{l+1}$};
 \draw[idi arrow] (h)--(f); \draw[idi arrow] (f)--(add); \draw[idi arrow] (add)--(out);
 \draw[idi arrow] (h.north)--++(0,0.8) -| node[pos=0.25,above,font=\small]{恒等捷径} (add.north);
\end{tikzpicture}
\caption{同尺寸残差块的计算路径。捷径传递输入，分支学习修正量，二者逐元素相加。图中省略了具体实现可能在相加后使用的激活。}
\label{fig:vision-residual-block}
\end{figure}

\paragraph{残差梯度不再只有单一路径。}
设各层$h_l\in\R^d$，且残差分支$F_l:\R^d\to\R^d$可微。令$J_l=\partial F_l/\partial h_l^{\mathsf T}$，则
\begin{align}
 \frac{\partial h_{l+1}}{\partial h_l^{\mathsf T}}&=I+J_l,\\
 \nabla_{h_l}\mathcal L&=(I+J_l)^{\mathsf T}\nabla_{h_{l+1}}\mathcal L.
\end{align}
连续两个块的前向Jacobian为$(I+J_{l+1})(I+J_l)=I+J_l+J_{l+1}+J_{l+1}J_l$，包含不经过任何残差分支的恒等项、经过一个分支的项和经过两个分支的项。相比普通网络只有$J_{l+1}J_l$，残差结构允许梯度走较短路径。若$\|J_l\|_2\leq\epsilon<1$，则对任意向量$v$有
\begin{equation}
 (1-\epsilon)\|v\|_2\leq\|(I+J_l)v\|_2\leq(1+\epsilon)\|v\|_2.
\end{equation}
小残差分支使单块接近恒等，但多块的上下界仍相乘，所以残差不是绝不消失或爆炸的定理。改变通道或空间尺寸时，要把捷径替换为投影$S_lh_l$，导数变为$S_l+J_l$，此时也不再有字面上的恒等路径。

\paragraph{相加要求尺寸一致。}
若残差输入和分支输出都是$8\times8\times16$，逐元素相加后仍是$8\times8\times16$。若残差分支改变空间尺寸或通道数，捷径也要用步幅卷积等投影变换得到相同尺寸。后文 U-Net 将使用通道拼接，同样两组特征拼接后为$8\times8\times32$，两种连接的输出维度和计算含义不同。

\subsection{从基本块到瓶颈块与轻量卷积}
ResNet-18 和 ResNet-34 主要使用由两个$3\times3$卷积构成的基本残差块。ResNet-50 等更深的版本使用瓶颈块：先以$1\times1$卷积减少通道，在较窄的表示上做$3\times3$卷积，再以$1\times1$卷积扩展通道。$1\times1$卷积不读取相邻空间位置，但会混合当前位置的不同通道。瓶颈设计把较昂贵的空间卷积放在较少通道上，使加深网络的计算代价更可控。残差块内部仍包含非线性变换，不能把整块等同为三个矩阵的单次线性乘法。

若相机侧计算资源更紧，MobileNet 一类轻量网络进一步使用深度可分离卷积。它先逐通道做空间卷积，再用$1\times1$卷积混合通道。忽略偏置，标准卷积每个位置需要约$k^2C_{\rm in}C_{\rm out}$次乘加，拆分后为$k^2C_{\rm in}+C_{\rm in}C_{\rm out}$，二者比值为$1/C_{\rm out}+1/k^2$。节省来自对空间与通道混合方式的限制，所以表达能力也发生改变。实际延迟还要在目标设备上测量，乘加数较少不保证同比例加速。

ResNet 解决深层表示的优化问题，瓶颈块和轻量卷积控制计算成本。它们都没有取消局部卷积的约束。接下来要问：若判断一个局部缺陷必须参考远处的部件，能否让区域之间直接交换信息？

\section{ViT：用注意力联系不同图像区域}
\subsection{从图像块到类别预测}
卷积通过逐层组合扩大观察范围。视觉 Transformer 将图像划分为一组图像块（patch），把每个块编码为一个向量单元（token），再通过自注意力更新这些向量\cite{dosovitskiy2021vit}。它仍可用于整图分类，改变的是骨干内部组织空间信息的方法。

设输入为$H\times W\times C$，每块大小为$P\times P$，且$H,W$均可被$P$整除。共有$N=HW/P^2$个块。第$i$块展平为列向量$x_i\in\R^{P^2C}$，线性投影为
\begin{equation}
 z_i=Ex_i+e_i^{\rm pos},\qquad E\in\R^{d\times P^2C}.
 \label{eq:vision-patch-embedding}
\end{equation}
$z_i\in\R^d$为块表示，$e_i^{\rm pos}\in\R^d$为位置嵌入。没有位置信息时，普通自注意力本身不区分相同图像块在网格中的排列。投影权重在各块之间共享，位置嵌入则标明块来自哪里。

一幅$32\times32\times3$图像按$8\times8$划块，得到$16$个块，每块包含$192$个数。投影到$d=64$后，按行放置这些向量得到$Z\in\R^{16\times64}$。对单头注意力，令$W_Q,W_K,W_V\in\R^{64\times64}$，计算
\begin{equation}
 Q=ZW_Q,\quad K=ZW_K,\quad V=ZW_V,\qquad
 A=\softmax\left(\frac{QK^{\mathsf T}}{\sqrt{64}}\right),\quad O=AV.
 \label{eq:vision-attention}
\end{equation}
$Q,K,V$分别称为查询、键和值，$A\in\R^{16\times16}$沿每一行归一化。$A_{ij}$决定更新第$i$块时读取第$j$块值向量的权重，输出$O$仍为$16\times64$。每个块因此可以在同一层读取任意其他块的信息。权重由当前表示计算，不等于人工给定的空间邻接关系，也不自动表示缺陷的因果来源。

原始 ViT 在块序列前加入一个可学习的分类单元，此时本例共有$17$个位置。经过多层带残差连接的注意力与逐位置前馈网络后，取分类单元的最终表示送入分类头，并用图像标签计算交叉熵。分类单元通过与图像块交互汇总信息。也可以使用块表示的平均值作分类，但这是不同的汇聚选择。用于检测或分割时，则要保留图像块与原网格的对应关系，另外设计空间输出头。

\subsection{为什么引入窗口注意力与层级表示}
全局注意力显式形成$N\times N$的得分矩阵，其主要交互计算量为$O(N^2d)$。若高、宽都加倍而块大小不变，块数增加到四倍，注意力得分数量增加到十六倍。细缺陷需要较高分辨率，这就使全局交互的成本成为直接限制。

Swin Transformer 将注意力限制在局部窗口内，并在相邻层之间移动窗口划分。例如前一层位于窗口边缘的两个块，下一层可以进入同一窗口，跨窗口信息由此逐层传递。若每个窗口含$M\times M$个单元且可整分，则有$N/M^2$个窗口，交互计算量为
\begin{equation}
 O\left(\frac{N}{M^2}(M^2)^2d\right)=O(NM^2d).
\end{equation}
固定$M$时，这部分成本随$N$线性增加。代价是单层只能读取窗口内信息。Swin 还逐级合并相邻单元，得到不同空间尺度的特征，使检测与分割可以使用层级输出。合并降低成本，也可能损失局部细节。

卷积把局部性和位置间的权重共享直接写入结构。原始 ViT 对这些空间规律的约束较弱，通常更依赖预训练规模或合适的训练策略。实际比较应控制输入分辨率、预训练数据和微调预算。即使拥有全局注意力，输入缩放或块投影已经丢失的信息也不能被自动补回。块投影在$d<P^2C$时不可能保留所有输入变化，但这不意味着小于块的缺陷一定不可见，结果取决于投影保留了什么以及标签要求模型识别什么。

\section{目标检测：以 Faster R-CNN 解释框级定位}
\subsection{为什么先提议区域，再分类和修正位置}
分类器每图输出一组类别概率，检测器却需要输出数量不固定的一组目标。一张图可能没有缺陷，也可能有多条裂纹。若把分类器在所有可能的矩形上重复运行，不仅计算量大，相邻矩形还会反复读取相同像素。

Faster R-CNN 继承区域卷积神经网络（Region-based Convolutional Neural Network，R-CNN）的区域识别思路，把主要图像计算集中在共享骨干中。其两阶段结构先找值得检查的区域，再判断每个区域属于什么类别并修正边界。区域提议网络（Region Proposal Network，RPN）在共享特征上，为预置参考框预测“目标还是背景”以及坐标偏移，筛选后产生候选区域。第二阶段对候选区域提取固定大小的特征，输出更具体的类别与进一步修正的框。候选区域也可能只包含背景，因此第二阶段仍需背景类。

\begin{figure}[htbp]
\centering
\begin{tikzpicture}[node distance=0.6cm]
 \node[idi input] (image) {图像};
 \node[idi state,right=of image] (feat) {共享特征};
 \node[idi process,right=of feat] (rpn) {RPN\\候选框};
 \node[idi process,right=of rpn] (roi) {区域特征\\固定大小};
 \node[idi output,right=of roi] (head) {类别\\修正框};
 \draw[idi arrow] (image)--(feat); \draw[idi arrow] (feat)--(rpn);
 \draw[idi arrow] (rpn)--(roi); \draw[idi arrow] (roi)--(head);
 \draw[idi arrow] (feat.south)--++(0,-0.7) -| node[pos=0.25,below,font=\small]{按候选位置读取特征} (roi.south);
\end{tikzpicture}
\caption{Faster R-CNN 的两阶段计算。区域提议给出在哪里读取，共享特征提供读取内容，第二阶段进行区域分类与框修正。}
\label{fig:vision-faster-rcnn}
\end{figure}

图\ref{fig:vision-faster-rcnn}中的区域特征来自整图共享特征，不需要为每个候选重新运行整个骨干。感兴趣区域（Region of Interest，RoI）的大小各不相同，原始 Faster R-CNN 使用 RoI 池化把它们转成固定大小的特征。池化中的坐标量化会影响精细对齐，后文 Mask R-CNN 使用的 RoIAlign 则通过插值采样减轻这一问题。

\subsection{框怎样匹配标签并参与训练}
锚框（anchor）是在特征位置预置的参考矩形，不是真实标注。假设一幅图有$100$个候选、两类缺陷，采用包含背景类的简化检测头，类别得分为$100\times3$，类别无关的框偏移为$100\times4$。训练先把候选与标注匹配，得到“这是背景还是哪类缺陷”及正样本的坐标目标，再计算损失。推理没有真实框，只能解码偏移、筛选低分框和消除重复；最终框数因图像而异。不同检测器的输出格式可能不同，例如独立的目标存在分数或类别专属回归头，不能把这个示例当作所有YOLO版本的接口。

训练必须先把参考框或候选框与真实标注匹配，才能确定分类标签和坐标目标。常用的几何依据是交并比（Intersection over Union，IoU）。两个框$A,B$的交并比为
\begin{equation}
 \operatorname{IoU}(A,B)=\frac{|A\cap B|}{|A\cup B|}.
\end{equation}
这里$|\cdot|$表示面积，IoU 取值为$[0,1]$。匹配到目标的候选是正样本，背景候选是负样本，阈值之间的候选还可以按协议忽略。RPN 的类别目标是目标或背景，第二阶段则区分具体缺陷类别。下面用第二阶段的简化目标解释分类与定位如何一起学习，完整训练还包括 RPN 的提议损失。

\paragraph{从几何框到可训练的回归目标。}
设参考框中心及宽高为$(a_x,a_y,a_w,a_h)$，真实框为$(g_x,g_y,g_w,g_h)$，宽高均严格为正且使用同一像素坐标系。相同相对位置误差在大框中对应更大的像素位移。为减小回归目标对框尺寸的依赖，定义无量纲目标
\begin{equation}
 t_x^*=\frac{g_x-a_x}{a_w},\quad t_y^*=\frac{g_y-a_y}{a_h},\quad
 t_w^*=\log\frac{g_w}{a_w},\quad t_h^*=\log\frac{g_h}{a_h}.
\end{equation}
网络输出$t\in\R^4$后，以$b_x=a_x+a_wt_x,b_y=a_y+a_ht_y,b_w=a_we^{t_w},b_h=a_he^{t_h}$解码。中心相对化降低尺度依赖，对数宽高保证解码尺寸为正；但$\partial b_w/\partial t_w=b_w$，极端输出仍可能带来不稳定。

对匹配到真实框的正样本，常用残差$e=t-t^*$的分段损失
\begin{equation}
 \rho_\beta(e)=\begin{cases}e^2/(2\beta),&|e|<\beta,\\|e|-\beta/2,&|e|\geq\beta,\end{cases}
 \qquad \beta>0.
\end{equation}
在接点$e=\pm\beta$，两段函数值均为$\beta/2$，两侧导数也分别相等，所以此处分段不妨碍一阶求导。其导数在小误差处为$e/\beta$，大误差处为$\operatorname{sign}(e)$，兼顾局部平滑与异常残差的有界影响。完整损失可写为
\begin{equation}
 \mathcal L=\frac1{N_{\rm cls}}\sum_i\ell_{\rm cls}(p_i,y_i)
 +\frac{\lambda}{\max(1,N_+)}\sum_{i:y_i>0}\sum_{j\in\{x,y,w,h\}}\rho_\beta(t_{ij}-t_{ij}^*).
\end{equation}
这里$p_i$是候选$i$的类别概率向量，$y_i=0$表示背景、$y_i>0$表示缺陷类别；$N_{\rm cls}>0$为参与分类计分的候选数，$N_+$为正样本数，$\lambda\geq0$控制定位项权重。背景样本没有真实框回归目标；分类与定位的归一化方式必须报告，否则相同$\lambda$可能代表不同梯度比例。

IoU直接比较几何重叠。设两个边长为$a$的正方形仅水平错位$\delta\in[0,a)$，则交集为$a(a-\delta)$，并集为$a(a+\delta)$，所以$\operatorname{IoU}=(a-\delta)/(a+\delta)$，在$\delta=0$右侧的导数为$-2/a$。这说明相同一像素偏移对小框更加苛刻。不相交框的普通IoU为零，在保持不相交的区域内不给出靠近方向，因此实际定位还会结合参数回归或加入外接框、中心距离等项。

\subsection{推理怎样得到数量不固定的检测结果}
推理时没有真实框可以匹配。模型先根据参考框解码预测偏移，再按类别分数筛掉低分框。多个候选可能指向同一条裂纹，因此还需要消除重复结果。非极大值抑制（Non-Maximum Suppression，NMS）通常逐类进行：先保留最高分框，再删除与它的 IoU 超过阈值的其他框，重复处理剩余候选。最终保留的框数由图像和阈值共同决定。NMS 不保证保留框正确，两个紧邻的真实缺陷也可能互相抑制。

检测评价同样需要空间匹配。先按置信度排序，同类预测框在给定 IoU 阈值下与真实框一对一匹配。未匹配预测计为误报，未匹配真实框计为漏检。平均精度（Average Precision，AP）汇总精确率与召回率的关系，平均精度均值（mean Average Precision，mAP）再按指定类别或 IoU 阈值取平均。不同协议的 mAP 不能直接比较。小框对偏移敏感，应另外报告不同缺陷尺寸下的召回。

\subsection{从两阶段到单阶段：省去逐候选识别}
两阶段结构明确地区分区域提议和区域识别，但逐候选计算也有成本。YOLO（You Only Look Once）代表的单阶段检测思路，直接从整图特征的多个位置预测类别与坐标，不再为每个提议单独运行第二阶段识别头。它适合对吞吐量有要求的检测任务，具体速度仍取决于骨干、分辨率、版本和设备。

早期 YOLO 将图像划为网格，由网格单元预测框、目标置信度和类别。一个局部单元可负责的目标数量受输出设计限制，密集目标和尺度变化由此成为困难。后续版本改进了多尺度输出与预测头，部分版本使用锚框，另一些使用无锚框预测。理解共同思路后，应按具体版本核对输出和正样本分配，不能把早期网格机制当作所有版本的接口。

单阶段检测会面对大量背景位置。若许多背景已经容易识别，它们仍可能在总损失中占据较大比例。RetinaNet 用焦点损失（focal loss）降低易样本权重：
\begin{equation}
 \mathcal L_{\rm focal}=-\alpha_t(1-p_t)^\gamma\log p_t.
\end{equation}
二分类时，$p_t$是对真实标签的预测概率，即正样本取$p$，负样本取$1-p$。$\alpha_t$为类别权重，$\gamma\geq0$控制降权强度。例如$\gamma=2$时，$p_t=0.99$的易样本权重为$0.0001$，$p_t=0.5$的样本权重为$0.25$。这改变了不同难度样本对训练的相对影响，标注错误也可能因此获得较大权重。

\paragraph{焦点损失怎样改变梯度。}
对正类样本设logit为$z$、$p=\sigma(z)$，则$\ell=-\alpha(1-p)^\gamma\log p$。不能只把交叉熵梯度乘以$(1-p)^\gamma$，因为这个权重本身依赖参数。乘积法则给出
\begin{equation}
 \frac{d\ell}{dz}=\alpha\gamma p(1-p)^\gamma\log p-\alpha(1-p)^{\gamma+1}.
\end{equation}
令$p=1-\varepsilon$并利用$\log(1-\varepsilon)\approx-\varepsilon$，得梯度约为$-\alpha(\gamma+1)\varepsilon^{\gamma+1}$；普通交叉熵则约为$-\alpha\varepsilon$。所以易样本梯度以更高阶衰减，困难样本占据更大的相对权重。但标注错误往往也是困难样本，焦点损失不具备自动甄别错误标签的能力，也不能保证输出概率仍然校准。

\subsection{FPN：让不同大小的缺陷使用不同尺度的特征}
小裂纹需要较密的空间采样，大范围表面损伤则需要较广的上下文。只在最深层特征上检测，可能已经失去小缺陷的位置。特征金字塔网络（Feature Pyramid Network，FPN）在骨干不同层的特征之间建立自顶向下的信息路径\cite{lin2017fpn}。

设骨干特征$C_l$的分辨率随层级$l$增加而降低。先用$1\times1$卷积$L_l$统一通道数，再从最高层开始融合：
\begin{equation}
 \widetilde P_L=L_L(C_L),\qquad
 \widetilde P_l=L_l(C_l)+\operatorname{Up}(\widetilde P_{l+1}).
 \label{eq:vision-fpn}
\end{equation}
$\operatorname{Up}$将高层特征上采样到$C_l$的空间大小，两个分支通道相同后逐元素相加。通常再用$3\times3$卷积处理融合结果，得到检测使用的$P_l$。高层提供较广的语义信息，横向连接提供较密的空间信息。不同大小的目标可分配到不同层级上预测。FPN 可以接入两阶段或单阶段检测器，它改变特征来源，不独立规定最终输出格式。

上采样只是把数表放大。FPN 能利用的细节来自仍保留这些信息的横向分支，若缺陷在所有分支中都已被丢弃，融合不能把它真实恢复。这个限制同样存在于像素级分割中。

\section{图像分割：以 U-Net 解释像素级预测}
\subsection{为什么解码器需要编码器的细节}
矩形框可以圈出裂纹，却同时包含很多正常背景，不能直接用于测量裂纹面积。分割将预测目标改为与图像对应的像素标签。对细长裂纹，模型既要判断局部变化是否属于裂纹，又要保持其边界与连续性。

U-Net 是具有编码器与解码器的卷积网络\cite{ronneberger2015unet}。编码器逐级下采样，形成越来越粗的空间特征。解码器逐级上采样，将表示恢复到输出网格。只把最深层特征放大时，精细位置难以确定，因此 U-Net 还把编码器同尺度特征传给解码器。图\ref{fig:vision-unet}展示了这一计算关系。

\begin{figure}[htbp]
\centering
\begin{tikzpicture}[node distance=0.65cm and 2.0cm]
 \node[idi state] (e1) {编码特征\\$32\times32\times16$};
 \node[idi state,below=of e1] (e2) {编码特征\\$16\times16\times32$};
 \node[idi process,below=of e2] (b) {深层特征\\$8\times8\times64$};
 \node[idi process,right=of e2] (d2) {上采样、拼接、卷积\\$16\times16\times32$};
 \node[idi process,right=of e1] (d1) {上采样、拼接、卷积\\$32\times32\times16$};
 \node[idi output,above=of d1] (out) {像素概率\\$32\times32$};
 \draw[idi arrow] (e1)--(e2); \draw[idi arrow] (e2)--(b);
 \draw[idi arrow] (b.east)-|(d2.south); \draw[idi arrow] (d2)--(d1); \draw[idi arrow] (d1)--(out);
 \draw[idi arrow,dashed] (e2)--node[above,font=\small]{同尺度拼接}(d2);
 \draw[idi arrow,dashed] (e1)--node[above,font=\small]{同尺度拼接}(d1);
\end{tikzpicture}
\caption{采用尺寸保持卷积的简化 U-Net。实线表示下采样或解码计算，虚线把编码特征送到同尺度解码层。尺寸为教学设置，未逐层复现原论文配置。}
\label{fig:vision-unet}
\end{figure}

例如，把$8\times8\times64$深层特征上采样并调整为$16\times16\times32$，与编码器的$16\times16\times32$特征按通道拼接，得到$16\times16\times64$，再通过卷积产生$16\times16\times32$解码特征。拼接保留两组特征供后续学习使用，和 ResNet 的逐元素相加不同。图中为了保持尺寸使用了填充，原始 U-Net 的无填充卷积需要裁切编码特征后再拼接。

同尺度连接使解码器可以同时读取细节与上下文，但也可能带入背景纹理。它提供了使用信息的路径，是否正确使用仍由训练决定。

\subsection{从概率图到掩码与训练损失}
语义分割为每个像素判类别，同类的两条裂纹可以用同一个标签；实例分割还要给不同裂纹分配不同实例身份。以一张$32\times32$图像的二值分割为例，模型输出$32\times32$个logit，即尚未转成概率的实数得分，经sigmoid成为同形状概率图，训练与同形状零一标签比较。推理时才按验证集确定的阈值变成硬掩码。多类互斥分割通常输出$32\times32\times K$，在每个像素的$K$个类别上归一化；不能在整幅图的全部像素上做一次softmax。

四个像素的真实掩码为$(1,1,0,0)$，预测硬掩码为$(1,0,1,0)$，则交集为一、并集为三，IoU为$1/3$，Dice为$2/(2+2)=1/2$。训练中的软Dice用概率代替硬预测，以便求导；阈值化通常不参与这条训练梯度。U-Net在解码端把低分辨率特征放大，与同尺度编码特征拼接，再经卷积输出像素得分，拼接提供的细节仍必须由数据教会网络如何使用。

\paragraph{像素似然、Dice与稀疏前景。}
设一幅图含$n$个像素，二值标签$y_i\in\{0,1\}$，模型输出$p_i=\sigma(z_i)$，其中$\sigma(z)=1/(1+e^{-z})$为 sigmoid 函数。将每个像素的标签建模为伯努利分布，在给定图像后采用乘积形式得到似然。取负对数并按像素数平均，得到二值交叉熵（Binary Cross Entropy，BCE）：
\begin{align}
 p(y\mid X)&=\prod_{i=1}^n p_i^{y_i}(1-p_i)^{1-y_i},\\
 \mathcal L_{\rm BCE}&=-\frac1n\sum_i[y_i\log p_i+(1-y_i)\log(1-p_i)],\\
 \frac{\partial\mathcal L_{\rm BCE}}{\partial z_i}&=\frac{p_i-y_i}{n}.
\end{align}
这个似然分解不表示真实裂纹像素相互独立，而是输出头采用的近似。前景极少时，大量背景项会主导平均损失。Dice从集合重叠出发：硬掩码的$D=2|A\cap B|/(|A|+|B|)$，将预测指示函数换成概率便得到软 Dice 损失：
\begin{equation}
 \mathcal L_{\rm Dice}=1-\frac{2\sum_i p_iy_i+\epsilon}{\sum_i p_i+\sum_i y_i+\epsilon}.
 \label{eq:vision-soft-dice}
\end{equation}
其中$\epsilon>0$是固定的小平滑常数，用于避免空掩码时分母为零，求导时不随$p_i$变化。记$N=2\sum_i p_iy_i+\epsilon$、$D_0=\sum_i p_i+\sum_i y_i+\epsilon$，则
\begin{equation}
 \frac{\partial\mathcal L_{\rm Dice}}{\partial p_j}
 =\frac{N-2y_jD_0}{D_0^2},\qquad
 \frac{\partial\mathcal L_{\rm Dice}}{\partial z_j}
 =\frac{N-2y_jD_0}{D_0^2}p_j(1-p_j).
\end{equation}
一个像素的梯度依赖全图预测总量，体现其区域级归一化，而不是各像素独立计分。对无前景图，$N=\epsilon$，当预测面积较大时惩罚梯度可能很小；可结合 BCE 并明确空掩码的评价约定。Dice也不是概率校准目标，因此不应把优化后的$p_i$未经验证就解释为可靠风险概率。

硬掩码下，令$J=\operatorname{IoU}$，由$|A|+|B|=|A\cup B|+|A\cap B|$得到$D=2J/(1+J)$。同一对掩码的两个指标单调对应，但先对图像平均再作非线性变换不等价；软Dice与阈值后二值IoU更不能直接替换。边界和拓扑验收正是补充这些面积指标遗漏的信息。

\subsection{U-Net 的改进与 Mask R-CNN 的实例输出}
U-Net 的横向连接直接拼接浅层与深层特征，两者表达的内容可能差别较大。U-Net++ 在连接路径中增加卷积与中间融合，使解码器逐步组合不同深度的表示。它增加了可学习的处理步骤，也增加计算和存储成本，是否改善边界要用相同数据与评价协议检验。

另一类困难是需要较大上下文，却不希望继续下采样。DeepLab 系列使用空洞卷积，在卷积核采样点之间留出间隔。例如空洞率为二的$3\times3$核覆盖五像素宽的范围，仍只有九个空间采样点。在式\eqref{eq:vision-receptive-field}中，将增量改为$d_l(k_l-1)j_{l-1}$即可计算空洞率$d_l$的感受野。空洞空间金字塔池化（Atrous Spatial Pyramid Pooling，ASPP）并行使用不同空洞率，组合不同范围的上下文。稀疏采样不等于观察了覆盖范围内的全部细节，边界仍要依赖足够密的特征。

这些结构输出语义掩码，同类相邻目标仍可能连在一起。若任务要分别统计每个缺陷的面积，可以使用 Mask R-CNN\cite{he2017mask}。它在 Faster R-CNN 的区域识别结构中增加一个与分类和框回归并行的掩码分支。RoIAlign 在共享特征上按候选框位置插值采样，减轻坐标取整造成的错位。每个正候选的掩码分支预测固定网格的类别掩码，训练时通常只对真实类别对应的掩码计二值交叉熵。推理时按预测类别取掩码，再映射回该框在原图中的位置。

Mask R-CNN 因此把“这是哪个目标”和“目标内部哪些位置属于前景”结合起来。它依赖前面的检测结果，漏掉的候选不会被掩码分支自行补回，较粗的区域网格也会限制细边界。语义分割适合整体前景区域分析，实例分割适合需要目标身份的测量。两者的选择由输出要求决定。
\section{回到质检任务：模型能解决什么，怎样验证}
\subsection{细小缺陷与有限标注}
同一张质检图，分类可以判断是否复检，检测可以指出复检位置，分割可以提供缺陷形状。三种输出的区别也决定了错误应如何检查。若分类正确而检测漏掉细裂纹，应先看裂纹是否在缩放或下采样时消失，再检查候选分配与筛选。如果框已经正确而掩码边缘断裂，则应检查解码分辨率、像素标签与掩码阈值。换一个更深的骨干不一定解决这些不同来源的问题。

标注不足时，可以使用在其他图像上预训练的骨干。冻结骨干，只训练最终线性分类层，称为线性探测。例如骨干每图输出$128$维向量，三类分类头包含$3\times128+3=387$个参数。全量微调还更新骨干，使原有表示适应当前图像。冻结骨干的线性头可以作为分类基线，检测和分割则必须训练各自的空间任务头。若预训练没有保留所需纹理，仅更新分类头可能不足。

预训练还可以利用无人工类别标签的图像，例如学习不同增强视图之间的对应关系，或重建被遮住的图像块。这些目标如何形成可复用表示，将在预训练章展开。图文对齐与根据图像生成报告，则放到多模态章讨论。本章先把类别、坐标和像素预测的机制与边界建立清楚。

增强也要服从当前任务。几何变换必须同步作用于图像、框和掩码，离散类别掩码缩放通常采用最近邻插值。若裁剪删掉全部裂纹，图像级标签也要相应检查。若颜色是缺陷的依据，过强颜色扰动可能破坏标签。增强需要模拟允许出现的变化，不能随意要求模型忽略任务所需信息。

\subsection{用独立数据检查类别、位置和边界}
设生产线上有四台相机。可以用前三台相机的早期批次训练、后期批次验证，完整留出第四台相机测试设备迁移。来自同一工件的近重复帧应放在同一集合，避免模型在测试时再次看到几乎相同的图。阈值、输入分辨率与后处理规则在验证集确定，固定测试集只用于最终评价。若要区分相机变化与时间变化的影响，还需要分别固定其中一个因素进行测试。

分类除准确率外，还应报告缺陷召回和正常图误报。例如十个真实缺陷中找对八个，另误报两个，精确率为$8/(8+2)=0.8$，召回率为$8/10=0.8$。检测须先完成前述类别与空间匹配，再统计这些数量。分类热力图可以帮助检查模型是否依赖批次标记，却不能替代经过框标注训练与评价的检测结果。

分割应同时检查面积重叠与几何质量。四个像素的例子说明 IoU 和 Dice 衡量集合重叠，却不记录错误发生在哪里。对细裂纹，少量断裂可能只改变很少像素，却影响连通性和长度估计。因此要结合边界距离、连通性和具体测量目标。面积过滤等后处理也可能删除真实小缺陷，应纳入完整预测流程一起评价。标注边缘本身模糊时，可保留多名标注者的分歧，避免把标注不确定性全部归因于模型。

置信度也需要单独验证。在大量被报为$0.8$缺陷概率的图像中，实际缺陷比例接近$80\%$，才支持这一概率解释。新相机与光照变化可能破坏这种对应。对无法自动判定的图像可以转人工复核，同时报告自动覆盖率。例如一百张图自动处理八十张，其中两张错误，覆盖率为$80\%$，已处理样本错误率为$2/80=2.5\%$，另二十张仍占用复核时间。

最后在目标设备上测量包含预处理与后处理的延迟，并保留图像的相机、工件、批次和标注版本，以及预测所用模型与缩放规则。这些记录使漏检可以回到具体环节检查。独立测试回答新数据上是否有效，记录与复核帮助解释失败发生在哪里，二者都不能由训练损失下降代替。

以下材料按“卷积与残差分类—图像块注意力—框级检测—像素分割”的顺序对应本章，可在完成相关算例后阅读或实践。
\section*{辅助学习材料}
\begingroup
\emergencystretch=3em
\begin{itemize}
\item \textbf{Convolutional Neural Networks (CNNs / ConvNets)}（Stanford CS231n，英文笔记）。仅在需要复习时阅读卷积层、输出尺寸与参数共享，回到本章“卷积的索引、参数与等变性”检查步幅和填充。\href{https://cs231n.github.io/convolutional-networks/}{在线阅读}。
\item \textbf{ResNet网络讲解}（中文视频）。对照“残差梯度不再只有单一路径”，区分恒等捷径与改变尺寸的投影捷径，解释残差块为何比单纯加深网络更易优化。\href{https://www.bilibili.com/video/BV1T7411T7wa}{观看讲解}；中文博客\href{https://zhuanlan.zhihu.com/p/42706477}{《详解残差网络》}可作图文补充，其标题和残差块主题已由必应索引核对，但知乎原页访问受限，未核读全文。
\item \textbf{Vision Transformer网络讲解}（中文视频）。重点看图像如何切成 patch、投影为 token 并加入位置信息，再对照“为什么引入窗口注意力与层级表示”计算分辨率变化带来的开销；这是 ViT 的图像建模讲解，不是泛自注意力入门。\href{https://www.bilibili.com/video/BV1Jh411Y7WQ}{观看讲解}。
\item \textbf{区域卷积神经网络（R-CNN）系列}（《动手学深度学习》，中文教程）。先读其中的 Fast R-CNN 小节与 RoI pooling 示例，将共享整图特征、候选区域、分类和框回归对应到“框怎样匹配标签并参与训练”。Fast R-CNN 使用选择性搜索；随后对照 Faster R-CNN 小节，辨明后者才引入区域提议网络（RPN），不可混称。\href{https://zh.d2l.ai/chapter\_computer-vision/rcnn.html}{在线阅读}。
\item \textbf{YOLO系列网络讲解(V1\textasciitilde V3)}（中文视频）。沿单阶段预测到多尺度检测的变化，比较网格、锚框与类别／位置输出，对照本章“从两阶段到单阶段：省去逐候选识别”和 FPN 小节；不要把早期版本的机制直接套给所有新版 YOLO。\href{https://www.bilibili.com/video/BV1yi4y1g7ro}{观看讲解}。
\item \textbf{U-Net网络讲解}（中文视频）。沿编码器、解码器与同尺度特征拼接追踪分辨率，回到“为什么解码器需要编码器的细节”检查细裂纹边界为何需要低层特征；并与 ResNet 的相加捷径区分。\href{https://www.bilibili.com/video/BV1Vq4y127fB/}{观看讲解}。上述四个模型视频均可从\href{https://github.com/WZMIAOMIAO/deep-learning-for-image-processing}{作者课程目录}查到配套代码。
\item \textbf{TorchVision Object Detection Finetuning Tutorial}（PyTorch 官方，英文实践）。跟随自定义数据集、预训练 Mask R-CNN 微调与评价流程，核对框、标签和掩码格式，对应本章“U-Net 的改进与 Mask R-CNN 的实例输出”及最后的验证讨论；它是检测／实例分割实践，不替代前面的 Fast R-CNN 原理解说。\href{https://docs.pytorch.org/tutorials/intermediate/torchvision\_tutorial.html}{在线阅读}。
\end{itemize}
\endgroup
\paragraph{自测与参考答案。}
参考框的中心及宽高为$(10,10,4,2)$，真实框为$(12,9,8,2)$，求回归目标并解码核对。参考答案：$t^*=(1/2,-1/2,\log2,0)$；解码中心为$(10+4/2,10-2/2)=(12,9)$，宽高为$(4e^{\log2},2e^0)=(8,2)$。四个目标均无量纲，不能把对数宽度$\log2$直接当作像素宽度。若候选被匹配为背景，则不应给它随意指定这组定位目标。

\paragraph{本章小结。}
本章从同一幅质检图像的三种输出展开。分类将空间特征汇聚成类别判断，ResNet 通过残差路径改善深层卷积网络的优化，ViT 则通过图像块注意力建立区域之间的联系。两种骨干都需要任务头，才能形成所需预测。

当任务要求位置时，Faster R-CNN 用共享特征、区域提议和区域识别输出检测框。单阶段检测省去逐候选识别，多尺度特征则让不同大小的缺陷有合适的预测层级。当任务要求边界时，U-Net 通过解码器与同尺度连接输出像素标签，Mask R-CNN 进一步把像素掩码与每个目标的身份结合。框匹配、像素损失和最终评价各自回答不同问题，不能用一个训练分数代表全部任务效果。

找到裂纹后，还需要结合设备关系、有效规程和维修记录判断如何处理。后续章节将分别讨论这些信息。视觉模型提供的类别、坐标和掩码应保留拍摄条件、尺度和模型版本，使其他模型或人员可以核对它们的依据。
